% Book 4, Year 2, chapter 15 -- Threads and synchronisation.
% Provenance (never printed): sources/book4/BRIEFS.md, chapter 15.
\chapter{Threads and Synchronisation}\label{ch:b2:threads-and-synchronisation}

A bus has one seat left. A customer arrives at window~A, and ten milliseconds later another at window~B. Both clerks
run the same program:

\omcode{code/bachelor-2/15-threads-and-synchronisation/listings/last_seat.c}{21}{32}{One window's
sale}\label{lst:b2:threads-and-synchronisation:sell}

\omoutput{code/bachelor-2/15-threads-and-synchronisation/outputs/last_seat.txt}{Two tickets for seat
1.}\label{out:b2:threads-and-synchronisation:hook}

Each window checked that a seat was left, and one was, at the moment it checked. Between the check and the sale the
other window checked too. Neither line of the program is wrong; their \omterm{def:b2:threads-and-synchronisation:interleaving}{interleaving} is. This chapter runs several
\omterm{def:b2:threads-and-synchronisation:thread}{threads} in one program, counts the ways their steps can interleave, explores all of them with a small checker, and
builds the tools that rule the bad ones out: \omterm{def:b2:threads-and-synchronisation:lock}{locks}, \omterm{def:b2:threads-and-synchronisation:condvar}{condition variables} and \omterm{def:b2:threads-and-synchronisation:semaphore}{semaphores}, in C, C++ and Rust.

\begin{recall}
Book 2: two \omterm{def:b2:processes-and-system-calls:process}{processes} sharing a resource can produce a \emph{\omterm{def:b2:threads-and-synchronisation:race}{race condition}}, a result that depends on the order of
their steps; a \emph{\omterm{def:b2:threads-and-synchronisation:mutex-property}{critical section}} uses the shared data, \emph{\omterm{def:b2:threads-and-synchronisation:mutex-property}{mutual exclusion}} lets one \omterm{def:b2:processes-and-system-calls:process}{process} at a time be
in it, and a \emph{\omterm{def:b2:threads-and-synchronisation:lock}{lock}} provides it; a \emph{deadlock} is a set of \omterm{def:b2:processes-and-system-calls:process}{processes} each
waiting for a resource another holds, and needs four conditions (\omterm{def:b2:threads-and-synchronisation:mutex-property}{mutual exclusion}, hold and wait, no preemption,
circular wait). \cref{ch:b2:processes-and-system-calls}: a \omterm{def:b2:processes-and-system-calls:process}{process} has its own memory. \cref{ch:b2:rust-ownership}:
an exclusive \omterm{def:b2:rust-ownership:borrow}{borrow} excludes every other \omterm{def:b2:rust-ownership:borrow}{borrow}, so two \omterm{def:b2:threads-and-synchronisation:thread}{threads} cannot both change one value.
\end{recall}

\section{Threads}\label{sec:b2:threads-and-synchronisation:threads}

\begin{definition}[Thread]\label{def:b2:threads-and-synchronisation:thread}
A \emph{thread}\index{thread} is one sequence of execution inside a \omterm{def:b2:processes-and-system-calls:process}{process}: it has its own registers, program counter
and stack, and shares with the other threads of the \omterm{def:b2:processes-and-system-calls:process}{process} everything else, its memory, its open files, its
\omterm{def:b2:files-pipes-signals:signal}{signal handlers}. \code{pthread_create(&t, NULL, f, arg)} starts a thread that runs \code{f(arg)};
\code{pthread_join(t, NULL)} waits for it to end.
\end{definition}

Two \omterm{def:b2:processes-and-system-calls:process}{processes} made by \code{fork} share nothing they do not arrange to share; two \omterm{def:b2:threads-and-synchronisation:thread}{threads} share all their global
variables and the heap. That is what makes \omterm{def:b2:threads-and-synchronisation:thread}{threads} fast to communicate, and what made \code{seats_left} a single
variable seen by both windows. On a machine with several cores the \omterm{def:b2:threads-and-synchronisation:thread}{threads} run at the same time; on one core the
scheduler switches between them. Either way, a \omterm{def:b2:threads-and-synchronisation:thread}{thread}'s steps follow each other in its own order, and the steps of
different \omterm{def:b2:threads-and-synchronisation:thread}{threads} mix in an order that nobody chooses.

\begin{definition}[Interleaving]\label{def:b2:threads-and-synchronisation:interleaving}
An \emph{interleaving}\index{interleaving} (or \emph{schedule}) of \omterm{def:b2:threads-and-synchronisation:thread}{threads} is one sequence of all their steps that keeps
the steps of each \omterm{def:b2:threads-and-synchronisation:thread}{thread} in that \omterm{def:b2:threads-and-synchronisation:thread}{thread}'s order.
\end{definition}

\begin{proposition}[Number of interleavings]\label{prop:b2:threads-and-synchronisation:interleavings}
Two \omterm{def:b2:threads-and-synchronisation:thread}{threads} of $m$ and $n$ steps that never wait have $\binom{m+n}{m}$ \omterm{def:b2:threads-and-synchronisation:interleaving}{interleavings}.
\end{proposition}

\begin{proof}
An \omterm{def:b2:threads-and-synchronisation:interleaving}{interleaving} is a sequence of $m + n$ positions. Choosing which $m$ positions hold the first \omterm{def:b2:threads-and-synchronisation:thread}{thread}'s steps
determines it entirely, since each \omterm{def:b2:threads-and-synchronisation:thread}{thread}'s steps fill their positions in their own order; every choice gives an
\omterm{def:b2:threads-and-synchronisation:interleaving}{interleaving}. There are $\binom{m+n}{m}$ choices.
\end{proof}

Two \omterm{def:b2:threads-and-synchronisation:thread}{threads} of 10 steps have $\binom{20}{10} = 184\,756$ \omterm{def:b2:threads-and-synchronisation:interleaving}{interleavings}; of 20 steps, about $1.4 \times 10^{11}$. A
program can be run thousands of times and meet only a few of them: testing cannot show that a concurrent program is
correct, and the hook's bug may appear once a year. For a small model of a program, though, all of them can be tried.

\section{Data races}\label{sec:b2:threads-and-synchronisation:races}

\code{counter++} is not one step. The processor loads the counter into a register, adds 1, and stores it back.
\code{race.c} writes the three steps out explicitly (with C11 \omterm{def:b2:threads-and-synchronisation:atomic}{atomic} loads and stores, defined below, so that the compiler keeps them
and the program stays defined), runs $k$ \omterm{def:b2:threads-and-synchronisation:thread}{threads} that each add 1 a million times, and compares:

\omcode{code/bachelor-2/15-threads-and-synchronisation/listings/race.c}{14}{22}{Load, add, store: an increment that
can lose updates}

\omoutput{code/bachelor-2/15-threads-and-synchronisation/outputs/race.txt}{One run on the author's laptop: the racy
result, the share of increments lost, and the time of three versions.}\label{out:b2:threads-and-synchronisation:race}

\begin{omfigure}
\begin{tikzpicture}[font=\small, s/.style={draw, rounded corners=2pt, minimum height=6mm, inner sep=3pt,
  font=\scriptsize}]
\node[anchor=east] at (-0.2,0.8) {thread A};
\node[anchor=east] at (-0.2,-0.2) {thread B};
\node[s, fill=omDef!12] (a1) at (1.0,0.8) {\code{r = x}\quad (reads 5)};
\node[s, fill=omThm!12] (b1) at (3.4,-0.2) {\code{r = x}\quad (reads 5)};
\node[s, fill=omDef!12] (a2) at (5.8,0.8) {\code{x = r + 1}\quad (writes 6)};
\node[s, fill=omThm!12] (b2) at (8.4,-0.2) {\code{x = r + 1}\quad (writes 6)};
\draw[-{Stealth[length=2mm]}, gray] (-0.1,-0.85) -- (10.2,-0.85) node[right, font=\scriptsize] {time};
\node[font=\scriptsize, align=center] at (5.0,-1.4) {\code{x} was 5; after two increments it is 6: A's update is
lost, overwritten by B};
\end{tikzpicture}
\par\smallskip{\small A lost update. Both \omterm{def:b2:threads-and-synchronisation:thread}{threads} read the old value before either writes; the second write
overwrites the first.}
\end{omfigure}

Alone, the \omterm{def:b2:threads-and-synchronisation:thread}{thread} loses nothing. With two \omterm{def:b2:threads-and-synchronisation:thread}{threads}, nearly half the increments vanish: whenever the two \omterm{def:b2:threads-and-synchronisation:thread}{threads}
overlap, each overwrites the other's work. The precise numbers change from run to run; the losses do not go away.

\begin{omfigure}
\begin{tikzpicture}
\begin{axis}[omrecord, width=0.7\linewidth, height=5cm, ybar, bar width=14pt,
    xlabel={threads}, ylabel={increments lost (\%)}, ymin=0, ymax=100, xtick=data,
    symbolic x coords={1,2,4,8}, enlarge x limits=0.18, nodes near coords,
    nodes near coords style={font=\scriptsize}]
  \addplot[fill=omDef!40, draw=omDef] table[x=threads, y=lost]
    {code/bachelor-2/15-threads-and-synchronisation/figures/lost.dat};
\end{axis}
\end{tikzpicture}
\par\smallskip{\small The lost increments of \cref{out:b2:threads-and-synchronisation:race}. The more \omterm{def:b2:threads-and-synchronisation:thread}{threads}
overlap, the larger the share of increments that one \omterm{def:b2:threads-and-synchronisation:thread}{thread} overwrites.}
\end{omfigure}

\begin{remark}[Race condition, between threads]\label{def:b2:threads-and-synchronisation:race}\index{race condition}
Book 2's \omterm{def:b2:threads-and-synchronisation:race}{race condition}, between \omterm{def:b2:processes-and-system-calls:process}{processes}, is the same bug between \omterm{def:b2:threads-and-synchronisation:thread}{threads}: the result depends on the order in which
their steps interleave, and some of those orders give a wrong result. The lost increments of \code{race.c} are one.
\end{remark}

\begin{definition}[Data race]\label{def:b2:threads-and-synchronisation:data-race}
Two accesses to the same memory location form a \emph{data race}\index{data race} when they come from different
\omterm{def:b2:threads-and-synchronisation:thread}{threads}, at least one is a write, at least one is not \omterm{def:b2:threads-and-synchronisation:atomic}{atomic}, and nothing orders one before the other. In C and C++ a
program with a data race has undefined behaviour.
\end{definition}

\begin{definition}[Atomic operation]\label{def:b2:threads-and-synchronisation:atomic}
An operation is \emph{atomic}\index{atomic operation} when no other \omterm{def:b2:threads-and-synchronisation:thread}{thread} can observe it half done: it takes effect
as one indivisible step of every \omterm{def:b2:threads-and-synchronisation:interleaving}{interleaving}. \code{atomic_fetch_add(&x, 1)} reads, adds and writes as one step.
\end{definition}

Our loads and stores are \omterm{def:b2:threads-and-synchronisation:atomic}{atomic} one by one, so \code{race.c} has no \omterm{def:b2:threads-and-synchronisation:data-race}{data race}; it still loses updates, because the
three steps together are not \omterm{def:b2:threads-and-synchronisation:atomic}{atomic}. The hook's \code{seats_left}, a plain \code{int}, does have a \omterm{def:b2:threads-and-synchronisation:data-race}{data race}, and a
tool can see it: compiled with ThreadSanitizer, the program reports each pair of unordered accesses it witnesses.
(On this machine the tool needs address randomisation turned off, which \code{setarch -R} does; \code{trim.sh} keeps
the lines that name the race.)

\omoutput{code/bachelor-2/15-threads-and-synchronisation/outputs/tsan.txt}{The race found while it
happens.}

Such a tool reports only the races of the \omterm{def:b2:threads-and-synchronisation:interleaving}{interleaving} that ran. To see them all, we can explore. A
\emph{model} of a program lists each \omterm{def:b2:threads-and-synchronisation:thread}{thread}'s steps; \code{interleavings.py} runs every \omterm{def:b2:threads-and-synchronisation:interleaving}{interleaving} of them,
recursively, and counts those that end badly. Two states that are the same (same position in every \omterm{def:b2:threads-and-synchronisation:thread}{thread}, same
variables) have the same future, so the count of each is computed once and remembered:

\omcode{code/bachelor-2/15-threads-and-synchronisation/listings/interleavings.py}{30}{62}{Exploring every schedule
from one state}\label{lst:b2:threads-and-synchronisation:explore}

A step is a label, a \emph{guard} (the step may run only when it holds, as for a wait) and an action. The counter of
\cref{out:b2:threads-and-synchronisation:race} becomes two steps per \omterm{def:b2:threads-and-synchronisation:thread}{thread}:

\omcode{code/bachelor-2/15-threads-and-synchronisation/listings/interleavings.py}{82}{87}{The counter's model}

\omoutput{code/bachelor-2/15-threads-and-synchronisation/outputs/interleavings.txt}{Every schedule of four small
models.}\label{out:b2:threads-and-synchronisation:explorer}

Two \omterm{def:b2:threads-and-synchronisation:thread}{threads} of 3 and 4 steps give $\binom{7}{3} = 35$ schedules, as the proposition says. Of the counter's 6
schedules, 4 lose an update, and the explorer prints one. The two other models are the subject of the next section.

\section{Mutual exclusion}\label{sec:b2:threads-and-synchronisation:mutex}

\begin{remark}[Critical section, mutual exclusion, between threads]\label{def:b2:threads-and-synchronisation:mutex-property}\index{critical section}\index{mutual exclusion}
Book 2 defined these words between \omterm{def:b2:processes-and-system-calls:process}{processes}; between \omterm{def:b2:threads-and-synchronisation:thread}{threads} they mean the same. A \omterm{def:b2:threads-and-synchronisation:mutex-property}{critical section} is a part of a
\omterm{def:b2:threads-and-synchronisation:thread}{thread}'s code that must not run while another \omterm{def:b2:threads-and-synchronisation:thread}{thread} runs a \omterm{def:b2:threads-and-synchronisation:mutex-property}{critical section} on the same data, and \omterm{def:b2:threads-and-synchronisation:mutex-property}{critical sections} are
in \omterm{def:b2:threads-and-synchronisation:mutex-property}{mutual exclusion} when at most one \omterm{def:b2:threads-and-synchronisation:thread}{thread} at a time is inside any of them.
\end{remark}

\begin{remark}[Lock, between threads]\label{def:b2:threads-and-synchronisation:lock}\index{lock}\index{mutex}
Book 2's \omterm{def:b2:threads-and-synchronisation:lock}{lock}, which provides \omterm{def:b2:threads-and-synchronisation:mutex-property}{mutual exclusion}, is for \omterm{def:b2:threads-and-synchronisation:thread}{threads} an object called a \emph{mutex}, with two operations
around \omterm{def:b2:threads-and-synchronisation:mutex-property}{critical sections}, \emph{\omterm{def:b2:threads-and-synchronisation:lock}{lock}} and \emph{unlock}: at most one \omterm{def:b2:threads-and-synchronisation:thread}{thread} at a time is between the \omterm{def:b2:threads-and-synchronisation:lock}{lock} it called
and its unlock, and a \omterm{def:b2:threads-and-synchronisation:thread}{thread} that calls \omterm{def:b2:threads-and-synchronisation:lock}{lock} while another holds it waits.
\end{remark}

The naive \omterm{def:b2:threads-and-synchronisation:lock}{lock}, a shared flag tested and then set, fails for the hook's reason: both \omterm{def:b2:threads-and-synchronisation:thread}{threads} can test the flag before
either sets it. The explorer finds 24 bad schedules among 42 (\cref{out:b2:threads-and-synchronisation:explorer}).
Testing and setting must not be separable, or something else must break the tie. Peterson's algorithm breaks it with
one more variable, \code{turn}, and only loads and stores:

% ledger: b2:peterson-1981
\omcode{code/bachelor-2/15-threads-and-synchronisation/listings/locks.c}{8}{23}{Peterson's \omterm{def:b2:threads-and-synchronisation:lock}{lock} for \omterm{def:b2:threads-and-synchronisation:thread}{threads} 0 and
1}\label{lst:b2:threads-and-synchronisation:peterson}

A \omterm{def:b2:threads-and-synchronisation:thread}{thread} announces that it wants to enter, then gives the turn to the other, and waits while the other wants to enter
and has the turn. The explorer's model of it has no bad schedule. The proof holds for any number of passes:

\begin{theorem}[Peterson's algorithm]\label{thm:b2:threads-and-synchronisation:peterson}
Assume that every load and store of \code{flag} and \code{turn} is one step of a single \omterm{def:b2:threads-and-synchronisation:interleaving}{interleaving} of the two
\omterm{def:b2:threads-and-synchronisation:thread}{threads} (sequential consistency). Then (1) the two \omterm{def:b2:threads-and-synchronisation:thread}{threads} are never both in their \omterm{def:b2:threads-and-synchronisation:mutex-property}{critical sections}; (2) a \omterm{def:b2:threads-and-synchronisation:thread}{thread}
that waits enters before the other \omterm{def:b2:threads-and-synchronisation:thread}{thread} has entered twice more.
\end{theorem}

\begin{proof}
(1) Suppose both are in their \omterm{def:b2:threads-and-synchronisation:mutex-property}{critical sections} at some moment. Each has written \code{turn} on its way in; say
\omterm{def:b2:threads-and-synchronisation:thread}{thread} 0 wrote last, \code{turn = 1}, so \code{turn} is still 1 (a \omterm{def:b2:threads-and-synchronisation:thread}{thread} writes \code{turn} only on its way in). Look
at \omterm{def:b2:threads-and-synchronisation:thread}{thread} 0's last test of the wait condition, which came after its write of \code{turn}. \omterm{def:b2:threads-and-synchronisation:thread}{Thread} 1 had already set
\code{flag[1]} before writing \code{turn}, which happened before \omterm{def:b2:threads-and-synchronisation:thread}{thread} 0's write, and has not cleared it since, being
in its \omterm{def:b2:threads-and-synchronisation:mutex-property}{critical section}. So the test found \code{flag[1]} true and \code{turn == 1}, and \omterm{def:b2:threads-and-synchronisation:thread}{thread} 0 kept waiting:
contradiction. The case where \omterm{def:b2:threads-and-synchronisation:thread}{thread} 1 wrote last is symmetric.

(2) Let \omterm{def:b2:threads-and-synchronisation:thread}{thread} 0 wait. If \omterm{def:b2:threads-and-synchronisation:thread}{thread} 1 is not trying to enter, \code{flag[1]} is false, and \omterm{def:b2:threads-and-synchronisation:thread}{thread} 0's next test lets it
in. Otherwise \omterm{def:b2:threads-and-synchronisation:thread}{thread} 1 enters at most once more, and leaves by clearing \code{flag[1]}; to enter again it must write
\code{turn = 0}. From that write on, \omterm{def:b2:threads-and-synchronisation:thread}{thread} 0's wait condition, \code{flag[1]} and \code{turn == 1}, is false, so
\omterm{def:b2:threads-and-synchronisation:thread}{thread} 0 enters at its next test; and \omterm{def:b2:threads-and-synchronisation:thread}{thread} 1 waits, since \code{flag[0]} stays true and \code{turn} stays 0 until
\omterm{def:b2:threads-and-synchronisation:thread}{thread} 0 has entered and left.
\end{proof}

\begin{omfigure}
\begin{tikzpicture}[font=\scriptsize, b/.style={draw, rounded corners=2pt, fill=omDef!10, align=left, inner sep=3pt,
  text width=3.6cm}, v/.style={draw, fill=omThm!12, minimum width=1.3cm, minimum height=5mm},
  w/.style={-{Stealth[length=1.8mm]}}]
\node[b] (t0) at (0,0) {\textbf{thread 0}\\ 1. \code{flag[0] = true}\\ 2. \code{turn = 1}\\ 3. wait while
\code{flag[1]} and \code{turn == 1}\\ 4. critical section\\ 5. \code{flag[0] = false}};
\node[b] (t1) at (8.2,0) {\textbf{thread 1}\\ 1. \code{flag[1] = true}\\ 2. \code{turn = 0}\\ 3. wait while
\code{flag[0]} and \code{turn == 0}\\ 4. critical section\\ 5. \code{flag[1] = false}};
\node[v] (f0) at (4.1,0.8) {\code{flag[0]}}; \node[v] (tu) at (4.1,0) {\code{turn}};
\node[v] (f1) at (4.1,-0.8) {\code{flag[1]}};
\draw[w] (t0.east |- f0) -- node[above, pos=0.5] {writes} (f0.west);
\draw[w, dashed] (f0.east) -- node[above, pos=0.5] {reads} (t1.west |- f0);
\draw[w] (t1.west |- f1) -- node[below, pos=0.5] {writes} (f1.east);
\draw[w, dashed] (f1.west) -- node[below, pos=0.5] {reads} (t0.east |- f1);
\draw[{Stealth[length=1.8mm]}-{Stealth[length=1.8mm]}] (t0.east) -- (tu.west);
\draw[{Stealth[length=1.8mm]}-{Stealth[length=1.8mm]}] (t1.west) -- (tu.east);
\end{tikzpicture}
\par\smallskip{\small Peterson's protocol. Each \omterm{def:b2:threads-and-synchronisation:thread}{thread} writes its own flag, reads the other's, and both write and
read the shared \code{turn}. Whoever writes \code{turn} last waits, if the other also wants to enter.}
\end{omfigure}

For $n$ \omterm{def:b2:threads-and-synchronisation:thread}{threads}, Lamport's bakery algorithm works like the ticket machine at a counter: each \omterm{def:b2:threads-and-synchronisation:thread}{thread} takes a number
larger than every number it sees, and the smallest number is served first; ties, possible because two \omterm{def:b2:threads-and-synchronisation:thread}{threads} can
read the same numbers, go to the smaller \omterm{def:b2:threads-and-synchronisation:thread}{thread} index.

% ledger: b2:lamport-bakery-1974
\omcode{code/bachelor-2/15-threads-and-synchronisation/listings/locks.c}{34}{53}{The bakery: take a ticket, then
wait for every smaller one}\label{lst:b2:threads-and-synchronisation:bakery}

\begin{theorem}[Bakery algorithm]\label{thm:b2:threads-and-synchronisation:bakery}
Under sequential consistency, no two \omterm{def:b2:threads-and-synchronisation:thread}{threads} are ever in their \omterm{def:b2:threads-and-synchronisation:mutex-property}{critical sections} at once. Moreover, if \omterm{def:b2:threads-and-synchronisation:thread}{thread} $i$
finishes taking its ticket before \omterm{def:b2:threads-and-synchronisation:thread}{thread} $k$ starts taking its own, then $i$ enters its \omterm{def:b2:threads-and-synchronisation:mutex-property}{critical section} before $k$
(first come, first served).
\end{theorem}

\begin{proof}
Write $(a, i) < (b, k)$ for the order of pairs: $a < b$, or $a = b$ and $i < k$. We show: if $i$ is in its \omterm{def:b2:threads-and-synchronisation:mutex-property}{critical
section} while $k$ holds a ticket $n_k \ne 0$, then $(n_i, i) < (n_k, k)$. Two such \omterm{def:b2:threads-and-synchronisation:thread}{threads} would then each be before
the other, which is impossible.

When $i$ passed its loop for $k$, it first saw \code{choosing[k]} false at some moment $t$, then read
\code{number[k]}. If it read a ticket with $(n_k, k) > (n_i, i)$, that ticket is still $k$'s unless $k$ has left and
taken a new one since; a new ticket is taken after $n_i$ was written, so it is larger than $n_i$. If it read 0, $k$
had no ticket at that read; since $k$ was not choosing at $t$, its present ticket was taken by a choosing that began
after $t$, after $i$ wrote $n_i$, so $k$ read $n_i$ and took a larger number. In every case $(n_i, i) < (n_k, k)$.

For the second claim: $k$ starts choosing after $n_i$ is written, so it reads $n_i$ and takes $n_k > n_i$; while $i$
waits, $k$'s loop for $i$ waits too, since \code{number[i]} is nonzero and smaller.
\end{proof}

\omoutput{code/bachelor-2/15-threads-and-synchronisation/outputs/locks_demo.txt}{A plain \code{counter++} under each
\omterm{def:b2:threads-and-synchronisation:lock}{lock}: nothing is lost.}

\code{choosing} looks unnecessary, and is not: without it, a \omterm{def:b2:threads-and-synchronisation:thread}{thread} could read \code{number[k]} as 0 while $k$ is
still computing its ticket, enter, and then be overtaken by $k$'s small ticket (\cref{exo:b2:threads-and-synchronisation:11}).
Both algorithms spin, testing in a loop, and both rely on sequential consistency, which real processors give only to
\omterm{def:b2:threads-and-synchronisation:atomic}{atomic operations} of the strongest ordering; the Year 3 volume explains the weaker orders. In practice one uses the
system's \omterm{def:b2:threads-and-synchronisation:lock}{lock}, \code{pthread_mutex_t}, or a single \omterm{def:b2:threads-and-synchronisation:atomic}{atomic} instruction when the \omterm{def:b2:threads-and-synchronisation:mutex-property}{critical section} is one operation:

\omcode{code/bachelor-2/15-threads-and-synchronisation/listings/race.c}{24}{41}{The counter under a mutex, and as one
\omterm{def:b2:threads-and-synchronisation:atomic}{atomic} instruction}

Both are exact (\cref{out:b2:threads-and-synchronisation:race}), and both cost time: the racy loop is fast because the
\omterm{def:b2:threads-and-synchronisation:thread}{threads} do not cooperate. The mutex costs most, because a \omterm{def:b2:threads-and-synchronisation:thread}{thread} that finds it held is put to sleep by the kernel and
woken later.

\section{Waiting for a condition}\label{sec:b2:threads-and-synchronisation:conditions}

A \omterm{def:b2:threads-and-synchronisation:lock}{lock} keeps \omterm{def:b2:threads-and-synchronisation:thread}{threads} apart. Often a \omterm{def:b2:threads-and-synchronisation:thread}{thread} must also wait for something another \omterm{def:b2:threads-and-synchronisation:thread}{thread} will do: a consumer for an
item, a producer for a free slot.

\begin{definition}[Condition variable]\label{def:b2:threads-and-synchronisation:condvar}
A \emph{condition variable}\index{condition variable} lets a \omterm{def:b2:threads-and-synchronisation:thread}{thread} holding a \omterm{def:b2:threads-and-synchronisation:lock}{lock} wait for a condition on the data
the \omterm{def:b2:threads-and-synchronisation:lock}{lock} protects. \code{pthread_cond_wait(&c, &m)} releases \code{m} and puts the \omterm{def:b2:threads-and-synchronisation:thread}{thread} to sleep, in one step; when
another \omterm{def:b2:threads-and-synchronisation:thread}{thread} calls \code{pthread_cond_signal(&c)} (wake one waiter) or \code{pthread_cond_broadcast(&c)} (wake all),
the sleeper wakes and takes \code{m} again before returning.
\end{definition}

\begin{definition}[Bounded buffer]\label{def:b2:threads-and-synchronisation:bounded-buffer}
A \emph{bounded buffer}\index{bounded buffer} is a queue of at most $N$ items shared by \emph{producer} \omterm{def:b2:threads-and-synchronisation:thread}{threads}, which
add items, and \emph{consumer} \omterm{def:b2:threads-and-synchronisation:thread}{threads}, which remove them; a producer waits while the buffer is full, and a consumer
while it is empty.
\end{definition}

It is the standard case of waiting for a condition. Here the buffer is a ring of $N$ slots.

\omcode{code/bachelor-2/15-threads-and-synchronisation/listings/bounded_buffer.c}{14}{38}{Putting and taking under one
\omterm{def:b2:threads-and-synchronisation:lock}{lock}, with two conditions}

\omoutput{code/bachelor-2/15-threads-and-synchronisation/outputs/bounded_buffer.txt}{Two producers, two consumers,
200\,000 items through 8 slots.}

The condition is tested in a \code{while}, not an \code{if}: between the signal and the moment the woken \omterm{def:b2:threads-and-synchronisation:thread}{thread} takes
the \omterm{def:b2:threads-and-synchronisation:lock}{lock} again, another \omterm{def:b2:threads-and-synchronisation:thread}{thread} may have taken the item. The explorer finds such a schedule
(\cref{exo:b2:threads-and-synchronisation:6}).

% ledger: b2:ewd123-semaphores
\begin{definition}[Semaphore]\label{def:b2:threads-and-synchronisation:semaphore}
A \emph{semaphore}\index{semaphore} is a non-negative integer with two \omterm{def:b2:threads-and-synchronisation:atomic}{atomic operations}: $P$ (\code{sem_wait}) waits
until the value is positive and decreases it by 1; $V$ (\code{sem_post}) increases it by 1, waking a waiter if there
is one. A semaphore that starts at $N$ counts $N$ available resources.
\end{definition}

With two \omterm{def:b2:threads-and-synchronisation:semaphore}{semaphores} the ring needs no condition at all: \code{empty} counts the free slots and \code{full} the filled
ones, and the mutex only protects the indices.

\omcode{code/bachelor-2/15-threads-and-synchronisation/listings/semaphores.c}{15}{34}{The ring with \omterm{def:b2:threads-and-synchronisation:semaphore}{semaphores}: wait
for a slot, fill it, announce it}

\omoutput{code/bachelor-2/15-threads-and-synchronisation/outputs/semaphores.txt}{Every item arrives, and the
\omterm{def:b2:threads-and-synchronisation:semaphore}{semaphores} end as they started.}

\begin{proposition}[The buffer's invariant]\label{prop:b2:threads-and-synchronisation:semaphore}
Let \emph{produced} be the number of items written into the ring and \emph{consumed} the number read from it. At every
moment, $0 \le \text{produced} - \text{consumed} \le N$; so no slot is overwritten before it is read, and none is read
before it is written.
\end{proposition}

\begin{proof}
A \omterm{def:b2:threads-and-synchronisation:semaphore}{semaphore}'s value is its initial value plus its $V$s minus its completed $P$s, and is never negative. Each item read
was preceded by a completed $P$ on \code{full}, and each $V$ on \code{full} follows an item written:
$\text{consumed} \le P(\text{full}) \le V(\text{full}) \le \text{produced}$. Each item written was preceded by a $P$ on
\code{empty}, and each $V$ on \code{empty} follows an item read: $\text{produced} \le P(\text{empty}) \le N +
V(\text{empty}) \le N + \text{consumed}$. With the indices advancing modulo $N$, the writer is never more than $N$
slots ahead of the reader, nor behind it.
\end{proof}

\begin{omfigure}
\begin{tikzpicture}[font=\scriptsize]
\foreach \k in {0,...,7} {
  \pgfmathsetmacro{\a}{90 - 45*\k}
  \ifnum\k<3 \def\fillc{omDef!35}\else\ifnum\k>5 \def\fillc{omDef!35}\else\def\fillc{white}\fi\fi
  \fill[\fillc, draw=black] (\a-21:1.1) arc (\a-21:\a+21:1.1) -- (\a+21:1.9) arc (\a+21:\a-21:1.9) -- cycle;
  \node at (\a:1.5) {\k};
}
\draw[-{Stealth[length=2mm]}, omThm, thick] (90-45*3:2.7) node[right] {\code{in}: next to fill} -- (90-45*3:2.0);
\draw[-{Stealth[length=2mm]}, omThm, thick] (90-45*6:2.7) node[left] {\code{out}: next to empty} -- (90-45*6:2.0);
\node[align=left, anchor=west] at (3.6,0.4) {filled slots: 6, 7, 0, 1, 2\quad \code{full} = 5};
\node[align=left, anchor=west] at (3.6,-0.3) {free slots: 3, 4, 5\quad \code{empty} = 3};
\node[align=left, anchor=west] at (3.6,-1.0) {\code{empty} + \code{full} = 8 between operations};
\end{tikzpicture}
\par\smallskip{\small The ring of 8 slots. A producer takes a free slot (P on \code{empty}), fills slot
\code{in}, and announces it (V on \code{full}); a consumer does the reverse from \code{out}.}
\end{omfigure}

Two more patterns use the same tools.

\begin{definition}[Barrier]\label{def:b2:threads-and-synchronisation:barrier}
A \emph{barrier}\index{barrier} for $k$ \omterm{def:b2:threads-and-synchronisation:thread}{threads} has one operation, \emph{wait}, which returns in none of them until all
$k$ have called it; it is then ready for the next round.
\end{definition}

A \omterm{def:b2:threads-and-synchronisation:barrier}{barrier} separates work that proceeds in rounds; \code{pthread_barrier_wait} returns a special value to exactly one of
the $k$ \omterm{def:b2:threads-and-synchronisation:thread}{threads}:

\omcode{code/bachelor-2/15-threads-and-synchronisation/listings/barrier.c}{12}{27}{Rounds separated by two
\omterm{def:b2:threads-and-synchronisation:barrier}{barriers}}

\omoutput{code/bachelor-2/15-threads-and-synchronisation/outputs/barrier.txt}{Each round's report sees the four
parts of that round.}

\begin{definition}[Readers--writers lock]\label{def:b2:threads-and-synchronisation:rwlock}
A \emph{readers--writers lock}\index{readers--writers lock} can be held in two modes: by any number of \omterm{def:b2:threads-and-synchronisation:thread}{threads} at once
for reading, or by one \omterm{def:b2:threads-and-synchronisation:thread}{thread} alone for writing, never both.
\end{definition}

It can be built as a \omterm{def:b2:threads-and-synchronisation:lock}{lock} that counts readers. If readers keep arriving, a writer may wait forever.

\begin{definition}[Starvation]\label{def:b2:threads-and-synchronisation:starvation}
A \omterm{def:b2:threads-and-synchronisation:thread}{thread} \emph{starves}\index{starvation} when it waits forever for a resource that other \omterm{def:b2:threads-and-synchronisation:thread}{threads} keep obtaining,
although the program as a whole goes on.
\end{definition}

A writer-preferring \omterm{def:b2:threads-and-synchronisation:lock}{lock} stops admitting readers as soon as a writer waits (\cref{exo:b2:threads-and-synchronisation:9});
it can then starve readers instead. \omterm{def:b2:threads-and-synchronisation:starvation}{Starvation} differs from deadlock, where nobody goes on.

% ledger: b2:ewd310-philosophers
Five philosophers sit at a round table with a fork between each pair of plates; to eat, a philosopher needs both
neighbouring forks. If each takes the left fork and then the right, and all five sit down together, each holds one
fork and waits for the next: a deadlock. With timed \omterm{def:b2:threads-and-synchronisation:lock}{locks} the program reports it instead of hanging:

\omcode{code/bachelor-2/15-threads-and-synchronisation/listings/philosophers.c}{20}{37}{Left fork, then right, with a
time limit}

\omoutput{code/bachelor-2/15-threads-and-synchronisation/outputs/philosophers.txt}{The circle of
waits.}\label{out:b2:threads-and-synchronisation:deadlock}

\begin{proposition}[Lock order]\label{prop:b2:threads-and-synchronisation:order}
If the \omterm{def:b2:threads-and-synchronisation:lock}{locks} are numbered and every \omterm{def:b2:threads-and-synchronisation:thread}{thread} that holds several \omterm{def:b2:threads-and-synchronisation:lock}{locks} acquired them in increasing order, no set of
\omterm{def:b2:threads-and-synchronisation:thread}{threads} can be waiting for each other in a circle.
\end{proposition}

The proof is \cref{exo:b2:threads-and-synchronisation:12}. For the table, number the forks 0 to 4: every philosopher
takes the lower-numbered of their two forks first. Only philosopher 4, between forks 4 and 0, changes: fork 0 first.

\omoutput{code/bachelor-2/15-threads-and-synchronisation/outputs/philosophers_ordered.txt}{Ordered: every philosopher
eats, 100\,000 times.}

\begin{omfigure}
\begin{tikzpicture}[font=\scriptsize]
\draw[fill=omDef!6] (0,0) circle (1.9);
\foreach \i in {0,...,4} {
  \pgfmathsetmacro{\a}{90 - 72*\i}
  \pgfmathsetmacro{\f}{90 - 72*\i + 36}
  \draw[fill=white] (\a:1.25) circle (0.32); \node at (\a:1.25) {P\i};
  \node[draw, fill=omThm!20, rounded corners=1pt, inner sep=1.5pt] at (\f:1.45) {\i};
  \node at (\a:2.3) {};
}
\node[align=left, anchor=west, font=\scriptsize] at (2.6,0.9) {P$i$ sits between forks $i$ and $i + 1 \bmod 5$.};
\node[align=left, anchor=west, font=\scriptsize] at (2.6,0.2) {Left first: P0 holds 0 and waits for 1, P1 holds 1
and waits for 2,\\ \dots, P4 holds 4 and waits for 0: a circle.};
\node[align=left, anchor=west, font=\scriptsize] at (2.6,-0.85) {Lower number first: P4 takes 0 before 4, so it
waits for 0\\ while holding nothing, and P3 can take fork 4: the circle is broken.};
\end{tikzpicture}
\par\smallskip{\small The philosophers' table, forks numbered between the plates. The order of
\cref{prop:b2:threads-and-synchronisation:order} needs only one philosopher to change hands.}
\end{omfigure}

\begin{history}
% ledger: b2:ewd123-semaphores, b2:ewd123-period, b2:ewd310-period, b2:lamport-bakery-1974, b2:peterson-1981
Edsger Dijkstra's notes \emph{Cooperating sequential processes}, written in the mid-1960s, present Th.~J. Dekker's
solution to \omterm{def:b2:threads-and-synchronisation:mutex-property}{mutual exclusion} for two \omterm{def:b2:processes-and-system-calls:process}{processes}, the first correct one, and introduce the \omterm{def:b2:threads-and-synchronisation:semaphore}{semaphore} with its
operations $P$ and $V$. A few years later, his notes \emph{Hierarchical ordering of sequential processes} posed the
five dining philosophers. Leslie Lamport published the bakery algorithm in 1974; it needs no \omterm{def:b2:threads-and-synchronisation:atomic}{atomic operation} beyond
reading and writing a number. Gary Peterson's \omterm{def:b2:threads-and-synchronisation:thread}{two-thread} algorithm, far shorter than Dekker's, appeared in 1981 in a
two-page note titled \emph{Myths about the \omterm{def:b2:threads-and-synchronisation:mutex-property}{mutual exclusion} problem}.
\end{history}

\section{The same programs in C++ and Rust}\label{sec:b2:threads-and-synchronisation:languages}

C++ wraps the same objects in classes. A \code{std::lock_guard} \omterm{def:b2:threads-and-synchronisation:lock}{locks} in its \omterm{def:b2:cpp-classes-raii:constructor}{constructor} and unlocks in its
\omterm{def:b2:cpp-classes-raii:constructor}{destructor}, so no path out of the block, \code{return} or \omterm{def:b2:cpp-classes-raii:exception}{exception}, can keep the \omterm{def:b2:threads-and-synchronisation:lock}{lock} (\cref{ch:b2:cpp-classes-raii}):

\omcode{code/bachelor-2/15-threads-and-synchronisation/listings/counter_threads.cpp}{14}{27}{Four
\code{std::thread}s, a guarded counter and an \omterm{def:b2:threads-and-synchronisation:atomic}{atomic} one}

\omoutput{code/bachelor-2/15-threads-and-synchronisation/outputs/counter_threads.txt}{Both exact.}

% ledger: b2:cpp-counting-semaphore
\code{std::counting_semaphore} gives the \omterm{def:b2:threads-and-synchronisation:semaphore}{semaphore} ring as a \omterm{def:b2:cpp-templates:template}{class template}:

\omcode{code/bachelor-2/15-threads-and-synchronisation/listings/bounded_queue.hpp}{13}{35}{A bounded queue in C++20}

\omoutput{code/bachelor-2/15-threads-and-synchronisation/outputs/pc_cpp.txt}{Two producers and two consumers on a
\code{BoundedQueue<long, 8>}.}

C and C++ accept a \omterm{def:b2:threads-and-synchronisation:data-race}{data race} and leave its behaviour undefined. Rust refuses it at compile time. A \omterm{def:b2:threads-and-synchronisation:thread}{thread} started by
\code{thread::spawn} may outlive the function that started it, so it may not \omterm{def:b2:rust-ownership:borrow}{borrow} that function's variables:

\omcode{code/bachelor-2/15-threads-and-synchronisation/listings/race_refused.rs}{5}{14}{A \omterm{def:b2:threads-and-synchronisation:thread}{thread} that \omterm{def:b2:rust-ownership:borrow}{borrows}
\code{count}}

\omoutput[1-12]{code/bachelor-2/15-threads-and-synchronisation/outputs/race_refused.txt}{The first of two errors:
the \omterm{def:b2:threads-and-synchronisation:thread}{thread} could outlive \code{count}.}

The compiler suggests \code{move}; for an \code{i32}, that would give the \omterm{def:b2:threads-and-synchronisation:thread}{thread} its own copy, and \code{main} would
print 0, which is safe and not what we meant. Sharing a counter takes \omterm{def:b2:cpp-copy-move:shared-ptr}{shared ownership}, \code{Arc}, of a value that
can be changed only through a \omterm{def:b2:threads-and-synchronisation:lock}{lock}, \code{Mutex}:

\omcode{code/bachelor-2/15-threads-and-synchronisation/listings/arc_mutex.rs}{5}{20}{\code{Arc<Mutex<u64>>}: the
only way in is \code{lock()}}

\omoutput{code/bachelor-2/15-threads-and-synchronisation/outputs/arc_mutex.txt}{Exact.}

The data is inside the mutex, so it cannot be reached without the \omterm{def:b2:threads-and-synchronisation:lock}{lock}; the guard returned by \code{lock()} unlocks
when dropped. Rust also offers channels, queues between \omterm{def:b2:threads-and-synchronisation:thread}{threads}; a \code{sync_channel(8)} is the \omterm{def:b2:threads-and-synchronisation:bounded-buffer}{bounded buffer}, and
its receiving end sees the end of the stream when every sender has been dropped, like a \omterm{def:b2:files-pipes-signals:pipe}{pipe}'s reader
(\cref{ch:b2:files-pipes-signals}):

\omcode{code/bachelor-2/15-threads-and-synchronisation/listings/channel_pc.rs}{5}{18}{Producers and a consumer joined
by a channel}

\omoutput{code/bachelor-2/15-threads-and-synchronisation/outputs/channel_pc.txt}{All 200\,000 items.}

\begin{inpractice}[Ticket offices, banks and shared counters]
Every booking system meets the hook's race: two agents, one seat. Databases settle it with transactions
(\cref{ch:b2:database-design}), servers with \omterm{def:b2:threads-and-synchronisation:lock}{locks} or with one \omterm{def:b2:threads-and-synchronisation:atomic}{atomic operation} such as ``decrement if positive''. The
working rules are short: protect every shared variable by one \omterm{def:b2:threads-and-synchronisation:lock}{lock}, always the same; hold it briefly; acquire several
\omterm{def:b2:threads-and-synchronisation:lock}{locks} in one global order; wait for conditions in a \code{while}; and prefer a language or a tool that finds \omterm{def:b2:threads-and-synchronisation:data-race}{data
races} for you.
\end{inpractice}

\omimg{book4/ai/b2-threads-and-synchronisation-ticket-office.jpg}{0.62\linewidth}{A ticket office with three
windows: three clerks selling from one stock of seats.}

\section{Exercises}\label{sec:b2:threads-and-synchronisation:exercises}

\begin{exercise}[$\star$]\label{exo:b2:threads-and-synchronisation:1}
Two \omterm{def:b2:threads-and-synchronisation:thread}{threads} each run \code{r = x; x = r + 1;} twice, with \code{x} shared and starting at 0. Which final values of
\code{x} are possible? Give a schedule for each, and say how many schedules there are.
\end{exercise}

\begin{exercise}[$\star$]\label{exo:b2:threads-and-synchronisation:2}
This program stops after its second line. Find the bug and fix it in two ways.
\omcode{code/bachelor-2/15-threads-and-synchronisation/listings/exo_early_return.c}{9}{17}{A withdrawal that may be
refused}
\omoutput{code/bachelor-2/15-threads-and-synchronisation/outputs/exo_early_return.txt}{The third call never
returns.}
\end{exercise}

\begin{exercise}[$\star$]\label{exo:b2:threads-and-synchronisation:3}
Write a C++ class \code{SafeCounter} with \code{add(d)}, \code{take(d)} (subtract $d$ only if the counter holds at
least $d$, and say whether it did) and \code{get()}, safe to call from many \omterm{def:b2:threads-and-synchronisation:thread}{threads}, using \code{std::lock_guard}.
\end{exercise}

\begin{exercise}[$\star$]\label{exo:b2:threads-and-synchronisation:4}
How many \omterm{def:b2:threads-and-synchronisation:interleaving}{interleavings} do two \omterm{def:b2:threads-and-synchronisation:thread}{threads} of 5 steps have? Three \omterm{def:b2:threads-and-synchronisation:thread}{threads} of 2 steps each? Generalise the proposition to
$k$ \omterm{def:b2:threads-and-synchronisation:thread}{threads} of $n_1, \dots, n_k$ steps.
\end{exercise}

\begin{exercise}[$\star\star$]\label{exo:b2:threads-and-synchronisation:5}
\code{s} and \code{t} are \omterm{def:b2:threads-and-synchronisation:semaphore}{semaphores} starting at 0. What can this program print? List every output.
\omcode{code/bachelor-2/15-threads-and-synchronisation/listings/exo_sem.c}{9}{28}{Two \omterm{def:b2:threads-and-synchronisation:thread}{threads}, two
\omterm{def:b2:threads-and-synchronisation:semaphore}{semaphores}}
\end{exercise}

\begin{exercise}[$\star\star$]\label{exo:b2:threads-and-synchronisation:6}
In \code{bounded_buffer.c}, replace \code{while (count == 0)} by \code{if (count == 0)}. Model one producer of two
items and two consumers of one item each, with a \omterm{def:b2:threads-and-synchronisation:lock}{lock} and a condition, and let the explorer find a schedule in which
a consumer takes from an empty buffer.
\end{exercise}

\begin{exercise}[$\star\star$]\label{exo:b2:threads-and-synchronisation:7}
Write a \omterm{def:b2:threads-and-synchronisation:barrier}{barrier} for $n$ \omterm{def:b2:threads-and-synchronisation:thread}{threads} from a mutex and a \omterm{def:b2:threads-and-synchronisation:condvar}{condition variable}, \code{barrier2_wait(b)}, that can be used
again and again, and returns 1 to exactly one \omterm{def:b2:threads-and-synchronisation:thread}{thread} per use. Why is the test \code{while (waiting < n)} wrong?
\end{exercise}

\begin{exercise}[$\star\star$]\label{exo:b2:threads-and-synchronisation:8}
Prove part (1) of \cref{thm:b2:threads-and-synchronisation:peterson} again, this time when the two lines of the entry
are swapped: \code{turn = j} first, then \code{flag[i] = true}. If the proof fails, give a schedule that breaks \omterm{def:b2:threads-and-synchronisation:mutex-property}{mutual
exclusion}.
\end{exercise}

\begin{exercise}[$\star\star$]\label{exo:b2:threads-and-synchronisation:9}
Write in Rust a \omterm{def:b2:threads-and-synchronisation:rwlock}{readers--writers lock} with writer preference, \code{RwLockW}, with \code{read_lock},
\code{read_unlock}, \code{write_lock} and \code{write_unlock}, from a \code{Mutex} and a \code{Condvar}. Which \omterm{def:b2:threads-and-synchronisation:thread}{threads}
can starve?
\end{exercise}

\begin{exercise}[$\star\star\star$]\label{exo:b2:threads-and-synchronisation:10}
The naive flag \omterm{def:b2:threads-and-synchronisation:lock}{lock} waits until \code{flag == 0}, then sets \code{flag = 1}. Exhibit, by hand, a schedule in which
both \omterm{def:b2:threads-and-synchronisation:thread}{threads} are in their \omterm{def:b2:threads-and-synchronisation:mutex-property}{critical sections}; then explain why the explorer finds 24 bad schedules out of 42, by
counting the schedules in which each \omterm{def:b2:threads-and-synchronisation:thread}{thread} passes its wait before the other sets the flag.
\end{exercise}

\begin{exercise}[$\star\star\star$]\label{exo:b2:threads-and-synchronisation:11}
Model the bakery for three \omterm{def:b2:threads-and-synchronisation:thread}{threads}, reading the other tickets one at a time, and check \omterm{def:b2:threads-and-synchronisation:mutex-property}{mutual exclusion} over every
schedule, for one and for two passes per \omterm{def:b2:threads-and-synchronisation:thread}{thread}. Then remove \code{choosing} and find a violation.
\end{exercise}

\begin{exercise}[$\star\star\star$]\label{exo:b2:threads-and-synchronisation:12}
Prove \cref{prop:b2:threads-and-synchronisation:order}. Then show that the philosophers' circle is broken as soon as
one philosopher takes the forks in the other order, whoever it is.
\end{exercise}

\section{Problem: The bus-ticket office with three windows}\label{sec:b2:threads-and-synchronisation:problem}

\begin{problem}[{Weekend problem --- The bus-ticket office with three windows}]\label{pb:b2:threads-and-synchronisation:1}
A bus company sells seats from three windows that share one count of seats left. The problem follows the hook's
race from its cause to a protocol checked under every \omterm{def:b2:threads-and-synchronisation:interleaving}{interleaving}, using \code{office_model.py}:

\omcode{code/bachelor-2/15-threads-and-synchronisation/listings/office_model.py}{10}{40}{Three protocols for one
customer's purchase}

\omoutput{code/bachelor-2/15-threads-and-synchronisation/outputs/office_model.txt}{Every schedule of each protocol,
for two windows.}\label{out:b2:threads-and-synchronisation:office}

\needspace{4\baselineskip}\medskip\noindent\textbf{Part I --- The race.}\par\nopagebreak
\begin{enumerate}
\item In \cref{lst:b2:threads-and-synchronisation:sell}, which lines form the check and which the sale? What does
\code{pause_ms(50)} stand for?
\item Window B's customer arrives 10 ms after A's. Explain the output of \cref{out:b2:threads-and-synchronisation:hook}
step by step. Why is \code{seats_left} 0 and not $-1$?
\item What would the program print if B arrived 60 ms after A?
\item The model \code{naive} has 3 steps per customer. Check that two windows with one customer each have 20
schedules.
\item Of these 20, 12 sell seat 1 twice. Count the 8 good ones directly.
\item Why does ThreadSanitizer report the race even in a run that happens to sell the seat once?
\end{enumerate}

\needspace{4\baselineskip}\medskip\noindent\textbf{Part II --- Locks.}\par\nopagebreak
\begin{enumerate}[resume, beginpenalty=10000]
\item In \code{locked}, the \omterm{def:b2:threads-and-synchronisation:lock}{lock} step has a guard. What does it model, and why is it one step?
\item Explain why the locked version with one customer per window has only 2 schedules.
\item With 4 customers per window and 6 seats, it has 70 schedules: explain $\binom{8}{4}$.
\item Rewrite \code{sell} of \code{last_seat.c} with a \code{pthread_mutex_t}. Which lines go inside?
\item With three windows, could Peterson's \omterm{def:b2:threads-and-synchronisation:lock}{lock} be used? What would you use instead?
\item A clerk holds the \omterm{def:b2:threads-and-synchronisation:lock}{lock} while printing the ticket (50 ms). What does that cost three windows, and how would you
shorten the \omterm{def:b2:threads-and-synchronisation:mutex-property}{critical section}?
\end{enumerate}

\needspace{4\baselineskip}\medskip\noindent\textbf{Part III --- A bounded queue.}\par\nopagebreak
\begin{enumerate}[resume, beginpenalty=10000]
\item The customers now queue in a hall that holds 8 people; a door lets someone in only if there is room. Which
\omterm{def:b2:threads-and-synchronisation:semaphore}{semaphore} is the door, and what is its initial value?
\item The three windows take customers from the hall. Which \omterm{def:b2:threads-and-synchronisation:semaphore}{semaphore} do they wait on?
\item Use \cref{prop:b2:threads-and-synchronisation:semaphore} to show the hall never holds more than 8 people.
\item Why is a mutex still needed around the ring's indices when two windows take at once?
\item A window must call two clerks' supervisors, holding \omterm{def:b2:threads-and-synchronisation:lock}{lock} $a$ then $b$, while another holds $b$ then $a$. What
can happen, and how does \cref{prop:b2:threads-and-synchronisation:order} prevent it?
\item In \code{semaphores.c}, what would happen if \code{put} took the mutex before \code{sem_wait(&empty)}?
\end{enumerate}

\needspace{4\baselineskip}\medskip\noindent\textbf{Part IV --- Every interleaving checked.}\par\nopagebreak
\begin{enumerate}[resume, beginpenalty=10000]
\item In the protocol \code{atomic}, the seat is taken by one indivisible step. Which C function would implement it?
\item Each window serves 4 customers in 2 steps each. Why does no guard reduce the number of schedules?
\item How many schedules are there, by \cref{prop:b2:threads-and-synchronisation:interleavings}?
\item What does the check \code{correct} require, and why is it enough for the company?
\item Why may the explorer remember a state's result and reuse it? Does that change the count?
\item State the result of the last line of \cref{out:b2:threads-and-synchronisation:office}.
\end{enumerate}
\end{problem}
